\subsection{Defense Repository Security Analysis}
\label{app:repository-risks}


The trusted computing base of CAITLYN includes the defense
repository, so an adversary who can tamper with the repository could
disable detection or install attacker-chosen rules.
This appendix analyzes the update path of the repository and the defenses
that keep it inside the trust boundary.

The local repository is written only by System~II and by an authenticated
operator.
System~II never exposes the raw trigger text to the generator, never
activates a candidate that failed deterministic verification, and never
promotes a shadow candidate without a clean observation window or explicit
approval.
The optional cloud-synchronization channel is off by default.
When enabled, a participating node can upload attack samples and
hand-written defense skills, but remote skills never auto-activate and
every contribution passes a human audit before cross-deployment release.
The channel is therefore a curated contribution pipeline rather than a
write channel into another host's repository.

Figure~\ref{fig:appendix-repository-guards} summarizes the six poisoning
defenses implemented on the local update path.
The residual risk is social and operational: a plausible-looking poisoned
contribution must still pass expert review, and an attacker who compromises
the operator's authentication can bypass the entire pipeline.
The adversarial model of Section~\ref{sec:AdaptiveAttacks} excludes both
paths by assumption, so an end-to-end poisoning campaign against the
curation pipeline remains an open research direction.
